\section{The harmonic comparison from Furusho's dyadic identities}
\label{sec:harmonic-dyadic-comparison}

Let $g=(1,f)$ be an element of full profinite $\mathrm{GT}$ with
$\rho_2^I(g)=0$.  We use the established equality between this
Ihara coordinate and the harmonic coordinate $\rho_2$, so the
pure harmonic conjugator $\beta\in\widehat F(a,b)$ has
abelianization zero.  Write $\theta(a)=b$, $\theta(b)=a$.
The fixed convention for this conjugator is
\begin{equation}\label{eq:harmonic-word-H}
f(a,b)=\theta(\beta)\beta^{-1}.
\end{equation}
Relation $(\mathrm{IV})$ holds on full GT by
\cite[Theorem~1]{H17}, and $(K_2)$ holds there by
\cite[Theorem~8.1]{Enriquez2007}. We apply these geometric
identities to the harmonic frame.

\subsection{The precise dyadic word identities}

The kernel of $\widehat F(x,y)\to\mathbb Z/2$ sending
$x\mapsto1$ and $y\mapsto0$ is free on $xyx^{-1},x^2,y$.  Since $f$ has zero
exponent sums, there is a unique word $f^+(W,X,Y)$ such that
\begin{equation}\label{eq:harmonic-word-1}
f(x,y)=f^+(xyx^{-1},x^2,y).
\end{equation}
Furusho's general definition includes the factor
$y^{-2\Psi_2^{(0)}(g)}$ on the left.  Here that factor is one,
because $\Psi_2^{(0)}=-\rho_2^I=0$.
The exact $(K_2)$ relation, in the form
\cite[\S3, equation~(2), p.~102]{Furusho2003}, gives
\begin{equation}\label{eq:harmonic-word-2}
f^+(1,a,b)=f(a,b).
\end{equation}
We use the following two identities in their full nonabelian
domains:
\begin{equation}\label{eq:harmonic-word-3}
f^+(W,X,Y)\,
 f^+(X^{-1}W^{-1}Y^{-1},Y,X)=1,
\end{equation}
and, whenever $WXYZ=1$,
\begin{equation}\label{eq:harmonic-word-4}
f^+(W,X,Y)f^+(Z,W,X)
 f^+(Y,Z,W)f^+(X,Y,Z)=1.
\end{equation}
These are respectively Furusho~\cite[Lemmas 4.1 and 4.2,
pp.~105--106]{Furusho2003}, specialized to
$\lambda=1$ and $\Psi_2^{(0)}=0$.
Their stated hypotheses are $(\mathrm{IV})$, and
$(\mathrm{I}), (\mathrm{II}), (\mathrm{IV})$, respectively.
They are not restricted to Galois elements.  Equation \eqref{eq:harmonic-word-3} is
an identity in the free profinite group on $W,X,Y$; equation
\eqref{eq:harmonic-word-4} is an identity in the free group on three of the four
variables with their indicated product relation.  We are not
using just their Fox derivatives or abelianizations.

\begin{lemma}\label{lem:dyadic-branch-word}
Put $h(a,b)=f^+(a,1,b)$.  Then
\[
 f(a,b)=h(a,b)h(b,a)^{-1}
 \quad\text{in }\widehat F(a,b).
\]
\end{lemma}
\begin{proof}
Work with $abc=1$.  Substitute $(W,X,Y,Z)=(1,b,c,a)$ in
\eqref{eq:harmonic-word-4}, obtaining
\begin{equation}\label{eq:harmonic-word-5}
f^+(1,b,c)f^+(a,1,b)f^+(c,a,1)f^+(b,c,a)=1.
\end{equation}
By \eqref{eq:harmonic-word-2}, the first factor is $f(b,c)$, and the second is
$h(a,b)$.  Applying \eqref{eq:harmonic-word-3} to $(W,X,Y)=(b,1,a)$ gives
\[
 f^+(b,1,a)f^+(b^{-1}a^{-1},a,1)=1.
\]
Since $b^{-1}a^{-1}=c$, the third factor of \eqref{eq:harmonic-word-5} is therefore
$h(b,a)^{-1}$.
Applying \eqref{eq:harmonic-word-3} to $(W,X,Y)=(b,c,a)$ gives
\[
 f^+(b,c,a)f^+(c^{-1}b^{-1}a^{-1},a,c)=1.
\]
Here $c^{-1}b^{-1}a^{-1}=1$, so \eqref{eq:harmonic-word-2} and the two-cycle
identity make the fourth factor of \eqref{eq:harmonic-word-5} equal to
$f(a,c)^{-1}=f(c,a)$.
Consequently
\[
 f(b,c)h(a,b)h(b,a)^{-1}f(c,a)=1,
\]
whence
\[
 h(a,b)h(b,a)^{-1}=f(b,c)^{-1}f(c,a)^{-1}.
\]
The normalized three-cycle identity is
$f(a,b)f(c,a)f(b,c)=1$ for $abc=1$.
It identifies the right side with $f(a,b)$, in precisely
the displayed order.  This proves the lemma.
\end{proof}

\begin{lemma}\label{lem:unique-swap-frame}
The map $\gamma\mapsto\theta(\gamma)\gamma^{-1}$ on
$\widehat F(a,b)$ is injective.  In particular
\[
 \beta(a,b)=h(b,a).
\]
\end{lemma}
\begin{proof}
If $\theta(\gamma_1)\gamma_1^{-1}=
\theta(\gamma_2)\gamma_2^{-1}$, write
$\gamma_2=\gamma_1 k$.  Cancellation gives $\theta(k)=k$.
The semidirect product
$\widehat F(a,b)\rtimes\langle s\rangle$, where $s^2=1$
and $sas=b$, is the free profinite product
$\widehat{\mathbb Z}\amalg C_2$, generated by $a,s$.
The centralizer of $s$ is the factor $\langle s\rangle$,
by the centralizer theorem for a nonidentity factor element
\cite[Theorem~B, p.~155]{HerfortRibes1985}.
Its intersection with $\widehat F(a,b)$ is trivial, so $k=1$.

By Lemma~\ref{lem:dyadic-branch-word}, the word
$\gamma(a,b)=h(b,a)$ satisfies
$\theta(\gamma)\gamma^{-1}=f(a,b)$.
Equation \eqref{eq:harmonic-word-H} and injectivity now give $\gamma=\beta$.
\end{proof}

\begin{proposition}\label{prop:normalized-GF1-from-dyadic}
For every $g$ as above, in the full signed profinite braid group
on generators $u,v$, with $s=uvu$, one has
\[
 \beta_g(u^2,v^2)^{-1}=f_g(u^2,s).
 \tag{GF1$_0$}
\]
\end{proposition}
\begin{proof}
Equation \eqref{eq:harmonic-word-2} shows that the $X$- and $Y$-exponent sums
of $f^+$ vanish, because $f$ is a commutator word.
Equation \eqref{eq:harmonic-word-1} then shows that its $W$-exponent sum also
vanishes. In particular its $X$-exponent is zero.

In the signed braid group, put $c=(uv)^3$.
The braid relation gives $su^2s^{-1}=v^2$ and $s^2=c$,
and $c$ is central. Specializing \eqref{eq:harmonic-word-1} at $(x,y)=(s,u^2)$
therefore gives
\[
 f(s,u^2)=f^+(v^2,c,u^2)=f^+(v^2,1,u^2)
         =h(v^2,u^2)=\beta(u^2,v^2).
\]
The second equality holds because a central argument can
be collected into its total exponent, which here is zero.
This assertion for profinite words follows by continuity
from the corresponding identity for finite words.
The two-cycle identity now gives
$f(u^2,s)=f(s,u^2)^{-1}=\beta(u^2,v^2)^{-1}$,
as required.
\end{proof}

\begin{corollary}\label{cor:harmonic-defect-zero}
On $D=\ker\lambda\cap\bigcap_q\ker\rho_q^I$, the finite
harmonic frame sign satisfies $w=\eta$.  Thus the previously
constructed relative character $\epsilon_h=w+\eta$ is zero
on the whole of full GT, and the remaining exceptional
character is $\chi_3=\epsilon_c=\tau+\eta$.
The cubic comparison in Section~\ref{sec:normalized-cubic-identity}
will also make this character zero.
\end{corollary}
\begin{proof}
Use the signed matrices over $\mathbb F_3$
\[
 U=\begin{pmatrix}1&1\\0&1\end{pmatrix},\qquad
 S=\begin{pmatrix}0&1\\-1&0\end{pmatrix}.
\]
The earlier finite comparison proves
$f_g(U,S)=(-I)^{\eta(g)}$ on $D$.
Let $J=\operatorname{diag}(-1,1)$.  Then
$JUJ^{-1}=U^{-1}=U^2$ and $JSJ^{-1}=-S$.
A commutator proword is unaffected by multiplying either
input by a central element: its exponent sum in that input
is zero.  Hence
\[
 f_g(U^2,S)=f_g(U^2,-S)
          =Jf_g(U,S)J^{-1}=(-I)^{\eta(g)}.
\]
Proposition~\ref{prop:normalized-GF1-from-dyadic} identifies
the left side with $q_3(\beta_g)^{-1}=(-I)^{w(g)}$.
This proves $w=\eta$ and $\epsilon_h|_D=0$.
Its established extension is zero on $G_{\mathbb Q}$, and
$\widehat{\mathrm{GT}}=D G_{\mathbb Q}$, so it is zero
everywhere.  The final equality follows from the already
proved decomposition $\chi_3=\epsilon_h+\epsilon_c$.
\end{proof}
